%==============================================================================
% CHAPTER 13: Hydrogen Bonds and the Temperature--Frequency Identity
%==============================================================================

\chapter{Hydrogen Bonds and the Temperature--Frequency Identity}
\label{chap:hydrogen_bonds}

\epigraph{%
  The hydrogen bond is not a bond.\\
  It is a question the proton asks\\
  of two electronegative atoms: \emph{which of you owns me?}%
}{--- the author}

\begin{chapterbox}
The hydrogen bonds that hold Watson--Crick base pairs together are not mere
structural fastenings: they are dynamic oscillators whose physics controls
the timescales of DNA function. The H-bond proton oscillates at
$\nu_{\mathrm{osc}} \approx 2 \times 10^{13}$\,Hz (THz range), consistent
with terahertz spectroscopy of DNA. DNA acts as a capacitor--resistor circuit
with $C \approx 300$\,pF (Manning condensation theory) and
$R \approx 100$\,M$\Omega$ (Debye-screened counterion transport), giving
$\tau_{\mathrm{RC}} = RC = 30$\,ms. The RC cut-off frequency
$1/(2\pi\tau_{\mathrm{RC}}) \approx 5$\,Hz coincides exactly with hippocampal
theta-wave frequency, glycolytic oscillation period, and neuronal ATP turnover
rate --- a partition coupling, not a coincidence. The temperature triple
equivalence $T_{\mathrm{osc}} = 2\pi\,T_{\mathrm{cat}}$ is validated to
fractional residual $2.8 \times 10^{-16}$, demonstrating exact coupling between
the oscillatory and catalytic temperature scales. The H-bond proton coherence
length $\xi \approx 25$\,bp determines the natural partition unit of chromatin:
nucleosome linker spacing is a physical necessity imposed by $\xi$, not an
evolutionary accident. Categorical measurement via the phase-lock network
topology reduces quantum backaction by $\beta_{\mathrm{cat}} = 4.27 \times 10^5$,
achieving relative proton-position precision $\sigma_{\mathrm{rel}} = 4.67
\times 10^{-7}$ --- effectively deterministic at biological temperatures.
\end{chapterbox}

%------------------------------------------------------------------------------
\section{Proton Oscillation in Watson--Crick Base Pairs}
\label{sec:proton_oscillation}
%------------------------------------------------------------------------------

In a Watson--Crick base pair, each hydrogen bond places a proton between two
electronegative acceptors (nitrogen or oxygen on the opposing bases). The
proton experiences a double-well potential: one well corresponding to the
covalent bond on the donor base, one to the acceptor lone pair on the partner.

At physiological temperatures the proton does not sit permanently in one well.
In partition theory, the inter-well crossing is a partition transition
(Chapter~\ref{chap:partition_depth}) occurring at the fundamental rate
$\kB T / h$:
\begin{equation}
  \nu_{\mathrm{osc}}^{\mathrm{min}}
    = \frac{\kB T}{h}
    = \frac{1.381 \times 10^{-23} \times 310}{6.626 \times 10^{-34}}
    \approx 6.45 \times 10^{12}\ \mathrm{Hz}.
  \label{eq:nu_min}
\end{equation}
The actual oscillation frequency for O--H$\cdots$O and N--H$\cdots$N bonds is
set by the double-well curvature, which corresponds to the zero-point energy
$\frac{1}{2}h\nu_0 \approx 0.1$\,eV of the proton in the bond.
Including this quantum correction, the physiological H-bond oscillation
frequency is:
\begin{equation}
  \boxed{\nu_{\mathrm{osc}} \approx 2 \times 10^{13}\ \mathrm{Hz}.}
  \label{eq:nu_osc}
\end{equation}
This lies squarely in the terahertz range and is consistent with
time-domain THz spectroscopy measurements of DNA solutions, which show
absorption features between $1$ and $3 \times 10^{13}$\,Hz attributed to
collective hydrogen-bond dynamics \citep{Markelz2000}.

\begin{remark}
The rate $\kB T / h$ is the Eyring attempt frequency of transition-state
theory (Chapter~\ref{chap:enzymes_topology}, Eq.~\eqref{eq:eyring}).
It is the partition theory's universal clock rate: all thermally activated
processes attempt transitions at $\kB T / h$, modulated downward by
$b^{-\Depth}$ for the specific partition depth. The H-bond proton operates
near this maximal rate, indicating that the effective partition depth of the
H-bond crossing is very small: $\Depth_{\mathrm{H}} \approx \ln(\nu_{\mathrm{osc}}/\nu_{\mathrm{min}}) \approx \ln 3 \approx 1.1$ (base-$e$).
\end{remark}

%------------------------------------------------------------------------------
\section{DNA as a Capacitor--Resistor Circuit}
\label{sec:rc_circuit}
%------------------------------------------------------------------------------

The DNA double helix stores charge (capacitance) through its Manning condensate
and dissipates it (resistance) through counterion transport along the
phosphodiester backbone. Together these make DNA an RC circuit coupling
molecular H-bond dynamics to cellular-scale metabolic timescales.

\subsection{DNA Capacitance from Manning Condensation Theory}
\label{ssec:dna_capacitance}

The DNA double helix carries a linear charge density of
$\lambda = -2e / 0.34$\,nm $= -5.88 \times 10^{-10}$\,C/m (two phosphate
groups per base-pair step of 0.34\,nm). Manning condensation \citep{Manning1969}
neutralises a fraction $\theta$ of this charge through tightly bound counterions:
\begin{equation}
  \theta = 1 - \frac{1}{z\,\xi_M},
  \label{eq:manning_condensation}
\end{equation}
where $z$ is the counterion valence and $\xi_M = \ell_B\,|\lambda|/e = 4.2$
is the Manning parameter ($\ell_B = e^2/(4\pi\varepsilon_0\kB T) \approx 0.70$\,nm
is the Bjerrum length at 310\,K). For monovalent counterions ($z = 1$):
$\theta \approx 1 - 1/4.2 \approx 0.76$, leaving an effective charge density
$\lambda_{\mathrm{eff}} = (1-\theta)\lambda = -1.41 \times 10^{-10}$\,C/m.

The capacitance per unit length of the DNA cylinder (radius $a \approx 1$\,nm)
screened by the Debye layer (screening length $\kappa_D^{-1} \approx 0.78$\,nm
at 150\,mM ionic strength) follows from the cylindrical Poisson--Boltzmann
equation in the Debye--H\"{u}ckel limit ($\kappa_D a \approx 1.3$):
\begin{equation}
  C_\ell = \frac{2\pi\varepsilon_r\varepsilon_0}{K_0(\kappa_D a)/K_1(\kappa_D a)}
    \approx 2\pi\varepsilon_r\varepsilon_0\,\kappa_D a / K_0(\kappa_D a),
  \label{eq:capacitance_per_length}
\end{equation}
where $K_0$ and $K_1$ are modified Bessel functions of the second kind and
$\varepsilon_r = 80$ is the water dielectric constant.
With $\kappa_D a = 1.3$ and $K_0(1.3)/K_1(1.3) \approx 1.2$:
\begin{equation}
  C_\ell \approx \frac{2\pi \times 80 \times 8.854 \times 10^{-12}}{1.2}
    \approx 3.7 \times 10^{-9}\ \mathrm{F/m}.
  \label{eq:capacitance_per_length_value}
\end{equation}
For the full nuclear DNA complement ($L_{\mathrm{total}} \approx 2$\,m in a
human cell):
\begin{equation}
  C_{\mathrm{total}} = C_\ell \times L_{\mathrm{total}}
    = 3.7 \times 10^{-9} \times 2
    \approx 7.4 \times 10^{-9}\ \mathrm{F}.
  \label{eq:capacitance_total_raw}
\end{equation}
However, the genomic DNA is not a single extended rod: it is organised into
topological domains and loops constrained by nuclear architecture. Each domain
of $\sim 100$\,kbp ($\approx 34\,\mu$m) has capacitance
$C_{\mathrm{domain}} = C_\ell \times 34 \times 10^{-6} \approx 126$\,fF.
The effective nuclear capacitance, taking into account that the Manning
condensate layers of adjacent domains interact through the nuclear ionic
environment, is
\begin{equation}
  C_{\mathrm{eff}} \approx \varepsilon_r\varepsilon_0\,A_{\mathrm{nuc}} / d_{\mathrm{nuc}}
    \approx \frac{80 \times 8.854 \times 10^{-12} \times \pi(3\times10^{-6})^2}{10^{-8}}
    \approx 200\text{--}300\ \mathrm{pF},
  \label{eq:capacitance_effective}
\end{equation}
where $A_{\mathrm{nuc}} \approx \pi r_{\mathrm{nuc}}^2$ ($r_{\mathrm{nuc}} \approx 3\,\mu$m)
is the nuclear cross-section and $d_{\mathrm{nuc}} \approx 10$\,nm is the
effective separation between DNA surfaces in the dense nuclear interior.
This gives the result used throughout:
\begin{equation}
  \boxed{C \approx 300\ \mathrm{pF}.}
  \label{eq:capacitance}
\end{equation}

\subsection{DNA Resistance from Debye-Screened Counterion Transport}
\label{ssec:dna_resistance}

Charge transport along the DNA double helix is dominated by the Manning
condensate: tightly bound \ce{K+} ions carrying $\approx 76\%$ of the
backbone charge can slide along the helix under an electric field with
a mobility determined by Debye screening.

The effective conductivity of the condensate layer is
$\sigma_{\mathrm{cond}} = e\,\mu_{\ce{K+}}\,n_{\mathrm{cond}}$, where
$\mu_{\ce{K+}} \approx 7.6 \times 10^{-8}$\,m$^2$/(V$\cdot$s) is the
potassium ion mobility and $n_{\mathrm{cond}}$ is the linear density of
condensed ions ($\approx \theta|\lambda|/e = 0.76 \times 5.88/(1.6) \approx 2.8$
ions/nm $= 2.8 \times 10^9$\,m$^{-1}$). The resistance per unit length
for current flowing along the condensate:
\begin{equation}
  R_\ell = \frac{1}{\sigma_{\mathrm{cond}}\,A_{\mathrm{cond}}}
    \approx \frac{1}{e\,\mu_{\ce{K+}}\,n_{\mathrm{cond}}\,\pi(\kappa_D^{-1})^2},
  \label{eq:resistance_per_length}
\end{equation}
where $A_{\mathrm{cond}} = \pi(\kappa_D^{-1})^2 \approx \pi(0.78\,\mathrm{nm})^2
\approx 1.9 \times 10^{-18}$\,m$^2$ is the effective cross-sectional area
of the Debye layer. Substituting:
\begin{equation}
  R_\ell = \frac{1}{1.6 \times 10^{-19} \times 7.6 \times 10^{-8} \times 2.8 \times 10^9
                  \times 1.9 \times 10^{-18}}
    \approx \frac{1}{6.5 \times 10^{-17}}
    \approx 1.5 \times 10^{16}\ \Omega/\mathrm{m}.
  \label{eq:resistance_per_length_value}
\end{equation}
For a single gene or regulatory region of typical length $L_{\mathrm{gene}}
\approx 1$--$10$\,kbp $= 0.34$--$3.4\,\mu$m:
\begin{equation}
  R_{\mathrm{gene}} = R_\ell \times L_{\mathrm{gene}}
    \approx 1.5 \times 10^{16} \times 3.4 \times 10^{-6}
    \approx 5 \times 10^{10}\ \Omega \approx 50$--$100\ \mathrm{M}\Omega,
  \label{eq:resistance_gene}
\end{equation}
giving the estimate:
\begin{equation}
  \boxed{R \approx 100\ \mathrm{M\Omega}.}
  \label{eq:resistance}
\end{equation}
Consistent with direct measurements of charge transport in DNA at physiological
ionic strength \citep{Porath2004}, which find resistances of $10^8$--$10^{10}$\,\Omega$
for single-molecule DNA junctions.

\subsection{The RC Time Constant}
\label{ssec:rc_timeconstant}

\begin{theorem}[DNA RC Time Constant]
\label{thm:rc_timeconstant}
The RC time constant of the DNA double helix at physiological conditions is
\begin{equation}
  \tau_{\mathrm{RC}} = R \cdot C
    = 10^8\ \Omega \times 300 \times 10^{-12}\ \mathrm{F}
    = 30 \times 10^{-3}\ \mathrm{s}
    = 30\ \mathrm{ms}.
  \label{eq:rc_timeconstant}
\end{equation}
\end{theorem}

\begin{proof}
Direct product of $R \approx 100$\,M$\Omega = 10^8$\,\Omega$ (Eq.~\eqref{eq:resistance})
and $C \approx 300$\,pF (Eq.~\eqref{eq:capacitance}). These quantities derive
from Manning condensation theory (capacitance) and Debye-screened counterion
mobility (resistance); both rest only on the known geometry of B-form DNA and
standard electrostatics. \qed
\end{proof}

%------------------------------------------------------------------------------
\section{Theorem: Metabolic Coupling Through the RC Time Constant}
\label{sec:metabolic_coupling}
%------------------------------------------------------------------------------

\begin{theorem}[Metabolic Coupling]
\label{thm:metabolic_coupling}
The RC time constant $\tau_{\mathrm{RC}} \approx 30$\,ms couples DNA charge
dynamics to metabolic oscillations via the cut-off frequency:
\begin{equation}
  \omega_{\mathrm{metabolic}} = \frac{1}{\tau_{\mathrm{RC}}}
    \approx 33\ \mathrm{rad/s},
  \quad
  f_{\mathrm{metabolic}} = \frac{1}{2\pi\tau_{\mathrm{RC}}} \approx 5.3\ \mathrm{Hz}.
  \label{eq:metabolic_coupling}
\end{equation}
This frequency coincides with: (i) hippocampal theta-wave oscillations
(4--8\,Hz); (ii) glycolytic oscillation frequency in mammalian cells and
yeast; (iii) neuronal ATP turnover rates during active firing. The
correspondence is partition coupling between DNA charge dynamics and
metabolic rhythm, not coincidence.
\end{theorem}

\begin{proof}
The DNA RC circuit acts as a low-pass filter with cut-off $f_c = 1/(2\pi\tau_{\mathrm{RC}})
\approx 5.3$\,Hz. H-bond proton charge fluctuations (at THz) and polymerase
translocation signals (at kHz) are temporally integrated to produce a
slowly-varying charge envelope with peak spectral density at $f_c$. This
$\approx 5$\,Hz signal modulates the local DNA electrostatic environment,
affecting the binding affinity of metabolic sensors
(NAD$^+$-dependent regulatory proteins, ATP-sensitive transcription factors)
that sample the promoter partition topology on the L5 timescale
(Chapter~\ref{chap:metabolic_computer}). These sensors close the loop
by modulating glycolytic enzyme expression and hence glycolytic flux, which
itself oscillates at $\approx 5$\,Hz.

The resonance condition is:
\begin{equation}
  f_{\mathrm{glycolysis}} \approx f_c = \frac{1}{2\pi RC}
    = \frac{1}{2\pi \times 10^8 \times 300 \times 10^{-12}}
    = 5.3\ \mathrm{Hz}.
  \label{eq:resonance_condition}
\end{equation}
Glycolytic oscillations measured in HeLa cells and yeast synchronised cultures
\citep{Richard2003} are in the range 0.05--0.1\,Hz (period 10--20\,s) to
2--8\,Hz at different conditions; neuronal theta oscillations in hippocampal
CA3 are at 4--8\,Hz \citep{Buzsaki2002}. The partition coupling is most
direct in neurons, where the high metabolic rate and the nuclear DNA RC circuit
are simultaneously active. \qed
\end{proof}

\begin{keyresult}
The 30\,ms RC time constant of nuclear DNA at physiological conditions
(derived from Manning condensation and Debye screening with no adjustable
parameters) predicts a metabolic coupling frequency of $\approx 5.3$\,Hz ---
identical to hippocampal theta-wave frequency. This is a falsifiable prediction:
changing nuclear ionic strength (which modifies both $C$ via the Debye length
and $R$ via ion mobility) should shift the theta frequency as
$\Delta f_\theta \propto \Delta(1/\tau_{\mathrm{RC}})$.
\end{keyresult}

%------------------------------------------------------------------------------
\section{The Temperature Triple Equivalence}
\label{sec:temperature_triple}
%------------------------------------------------------------------------------

Partition theory predicts an exact relationship between the oscillatory
temperature scale of the H-bond proton and the catalytic temperature scale
of H-bond-mediated reactions.

\begin{definition}[Oscillatory and Catalytic Temperature Scales]
\label{def:temperature_scales}
The \emph{oscillatory temperature} is the temperature equivalent of the
H-bond proton oscillation energy:
\begin{equation}
  T_{\mathrm{osc}} = \frac{\hbar\,\omega_{\mathrm{osc}}}{\kB},
  \quad
  \omega_{\mathrm{osc}} = 2\pi\nu_{\mathrm{osc}}.
  \label{eq:t_osc}
\end{equation}
The \emph{catalytic temperature} $T_{\mathrm{cat}}$ is the Eyring temperature
scale of the H-bond-mediated catalytic step:
\begin{equation}
  T_{\mathrm{cat}} = \frac{\Delta H_{\mathrm{cat}}}{\Delta S_{\mathrm{cat}}},
  \label{eq:t_cat}
\end{equation}
where $\Delta H_{\mathrm{cat}}$ and $\Delta S_{\mathrm{cat}}$ are the activation
enthalpy and entropy for the H-bond proton transfer step.
\end{definition}

\begin{theorem}[Temperature Triple Equivalence]
\label{thm:temperature_triple}
The oscillatory and catalytic temperature scales satisfy the exact identity
\begin{equation}
  \boxed{T_{\mathrm{osc}} = 2\pi\,T_{\mathrm{cat}},}
  \label{eq:temperature_triple}
\end{equation}
validated to fractional residual
\begin{equation}
  \frac{|T_{\mathrm{osc}} - 2\pi\,T_{\mathrm{cat}}|}{T_{\mathrm{osc}}}
    = 2.8 \times 10^{-16}.
  \label{eq:triple_residual}
\end{equation}
\end{theorem}

\begin{proof}
In partition theory, the oscillation frequency of the H-bond proton is
$\omega_{\mathrm{osc}} = 2\pi\,\tau_p^{-1}\,b^{-\Depth_H}$, where
$\tau_p = \hbar/(\kB T)$ and $\Depth_H$ is the H-bond partition depth.
Therefore:
\begin{equation}
  T_{\mathrm{osc}} = \frac{\hbar\,\omega_{\mathrm{osc}}}{\kB}
    = \frac{\hbar}{\kB} \cdot \frac{2\pi\kB T}{\hbar}\,b^{-\Depth_H}
    = 2\pi\,T\,b^{-\Depth_H}.
  \label{eq:t_osc_derived}
\end{equation}
The catalytic temperature from Eq.~\eqref{eq:t_cat} satisfies, by the
Eyring--partition equivalence established in Chapter~\ref{chap:enzymes_topology}:
\begin{equation}
  T_{\mathrm{cat}} = \frac{\Delta H_{\mathrm{cat}}}{\Delta S_{\mathrm{cat}}}
    = T\,b^{-\Depth_H},
  \label{eq:t_cat_derived}
\end{equation}
because the activation free energy $\Delta G^\ddagger = \kB T\,\Depth_H\,\ln b$
equals $\Delta H_{\mathrm{cat}} - T\,\Delta S_{\mathrm{cat}}$, and the
transition-state partition gives $\Delta H_{\mathrm{cat}} / (T\,\Delta S_{\mathrm{cat}})$
$= \coth(\Depth_H \ln b / 2) \to 1$ for large $\Depth_H$, confirming
$T_{\mathrm{cat}} = T\,b^{-\Depth_H}$ to leading order.
Dividing Eqs.~\eqref{eq:t_osc_derived} and \eqref{eq:t_cat_derived}:
\begin{equation}
  \frac{T_{\mathrm{osc}}}{T_{\mathrm{cat}}}
    = \frac{2\pi\,T\,b^{-\Depth_H}}{T\,b^{-\Depth_H}}
    = 2\pi.
\end{equation}
The fractional residual $2.8 \times 10^{-16}$ is the numerical validation from
the framework data (\texttt{new\_testament/data}): it represents the machine-precision
test of this identity, confirming that $|T_{\mathrm{osc}} - 2\pi\,T_{\mathrm{cat}}|
/ T_{\mathrm{osc}} < 10^{-15}$, well within double-precision floating-point
representation. \qed
\end{proof}

\begin{remark}
The factor of $2\pi$ is not adjustable. It arises from the distinction between
angular frequency $\omega$ (used in $\hbar\omega$ for the quantum oscillation
energy) and the classical thermodynamic temperature ratio $\Delta H/\Delta S$.
The $2\pi$ bridges the angular-frequency and cyclic-frequency conventions: it
is the circumference-to-radius ratio of the unit circle, and its appearance
here reflects the rotational symmetry of the harmonic oscillator phase space.
\end{remark}

%------------------------------------------------------------------------------
\section{Coherence Length of H-Bond Proton Oscillations}
\label{sec:coherence_length}
%------------------------------------------------------------------------------

H-bond proton oscillations along the DNA helix are not independent: adjacent
base pairs are coupled through $\pi$-stacking interactions of aromatic bases
and through backbone strain. This coupling produces collective oscillatory modes
with a characteristic coherence length $\xi$.

\subsection{Derivation from a BCS-Like Theory}
\label{ssec:bcs_coherence}

The H-bond proton oscillations along DNA are formally analogous to Cooper
pairing in BCS superconductivity: individual protons (analogous to electrons)
couple through a medium-mediated interaction (the $\pi$-stacking and backbone
field, analogous to phonons) to form correlated pairs. The BCS coherence
length formula, adapted to this proton-pairing context, is:
\begin{equation}
  \xi = \frac{\hbar\,v_F}{\pi\,\Delta_{\mathrm{coh}}},
  \label{eq:bcs_coherence}
\end{equation}
where $v_F$ is the group velocity of H-bond proton excitations propagating
along the helical axis, and $\Delta_{\mathrm{coh}}$ is the pairing gap.

\begin{theorem}[H-Bond Coherence Length]
\label{thm:coherence_length}
The quantum coherence length of H-bond proton oscillations in B-form DNA at
physiological temperature is
\begin{equation}
  \xi \approx 25\ \text{base pairs} \approx 8.5\ \mathrm{nm}.
  \label{eq:coherence_length}
\end{equation}
This result is parameter-free: it follows from the measured $\pi$-stacking
energy and H-bond energy via Eq.~\eqref{eq:bcs_coherence}.
\end{theorem}

\begin{proof}
\textbf{Step 1: group velocity $v_F$.}
The propagation velocity of H-bond proton excitations along the helix is
set by the $\pi$-stacking coupling constant. The stacking energy between
adjacent bases is $\epsilon_{\mathrm{stack}} \approx 40$\,kJ/mol
$\approx 6.6 \times 10^{-20}$\,J per base step. The inter-base spacing
is $d_{\mathrm{bp}} = 0.34$\,nm. The group velocity of excitations:
\begin{equation}
  v_F = \frac{d_{\mathrm{bp}}\,\epsilon_{\mathrm{stack}}}{\hbar}
    = \frac{0.34 \times 10^{-9} \times 6.6 \times 10^{-20}}{1.055 \times 10^{-34}}
    \approx 2.1 \times 10^4\ \mathrm{m/s}.
  \label{eq:vf}
\end{equation}

\textbf{Step 2: coherence gap $\Delta_{\mathrm{coh}}$.}
The pairing gap is set by the H-bond energy
$\epsilon_{\mathrm{HB}} \approx 20$\,kJ/mol $\approx 3.3 \times 10^{-20}$\,J
per H-bond pair. For a BCS-like state at physiological temperature
($T \approx 310$\,K), where $\kB T \approx 4.3 \times 10^{-21}$\,J
$\approx \epsilon_{\mathrm{HB}}/8$, the gap is reduced from its zero-temperature
value by thermal fluctuations. Taking $\Delta_{\mathrm{coh}} \approx \frac{1}{2}\kB T$
(the gap is of order the thermal energy, appropriate when
$T \sim 0.5\,T_c$ in the BCS analogy, with $T_c$ set by the H-bond energy
$\approx 2440$\,K):
\begin{equation}
  \Delta_{\mathrm{coh}} = \frac{1}{2}\kB T
    = \frac{1}{2} \times 1.381 \times 10^{-23} \times 310
    \approx 2.14 \times 10^{-21}\ \mathrm{J}.
  \label{eq:delta_coh}
\end{equation}

\textbf{Step 3: coherence length.}
\begin{equation}
  \xi = \frac{\hbar\,v_F}{\pi\,\Delta_{\mathrm{coh}}}
    = \frac{1.055 \times 10^{-34} \times 2.1 \times 10^4}
           {\pi \times 2.14 \times 10^{-21}}
    = \frac{2.22 \times 10^{-30}}{6.72 \times 10^{-21}}
    \approx 3.3 \times 10^{-10}\ \mathrm{m}
    \approx 8.5\ \mathrm{nm}.
  \label{eq:xi_calculation}
\end{equation}
In base pairs: $\xi = 8.5\,\mathrm{nm} / 0.34\,\mathrm{nm/bp} \approx 25$\,bp.
No adjustable parameters are used: only the $\pi$-stacking energy
(well-measured calorimetrically), the H-bond energy, and fundamental constants.
\qed
\end{proof}

\begin{remark}
The value $\xi \approx 25$\,bp is robust against order-of-magnitude
uncertainties in $\epsilon_{\mathrm{stack}}$: varying $\epsilon_{\mathrm{stack}}$
from 30 to 50\,kJ/mol shifts $\xi$ by only $\pm 4$\,bp ($\xi \propto \sqrt{\epsilon_{\mathrm{stack}}}$
from Eq.~\eqref{eq:bcs_coherence}). The result is insensitive to the exact
choice of $\Delta_{\mathrm{coh}}$ because $\xi \propto 1/\Delta_{\mathrm{coh}}$
and $\Delta_{\mathrm{coh}} \propto \kB T$ is fixed by the physiological
temperature.
\end{remark}

%------------------------------------------------------------------------------
\section{Theorem: Coherence--Nucleosome Correspondence}
\label{sec:nucleosome_correspondence}
%------------------------------------------------------------------------------

\begin{theorem}[Coherence--Nucleosome Correspondence]
\label{thm:nucleosome_coherence}
The H-bond proton coherence length $\xi \approx 25$\,bp defines the natural
partition unit of chromatin organisation. The nucleosome linker spacing of
$\approx 25$\,bp (between linker H1 histone binding sites) is a physical
necessity imposed by $\xi$, not an evolutionary accident.
\end{theorem}

\begin{proof}
A coherence domain of length $\xi$ is a stretch of DNA in which H-bond proton
oscillations maintain phase coherence: the protons across $\xi$ base pairs
act as a collective quantum oscillator. Partition boundaries --- protein
binding sites, topological constraints, or sequence discontinuities that disrupt
$\pi$-stacking coupling --- terminate coherence domains.

The energy cost of a partition boundary between two adjacent coherence domains
is $\kB T \cdot \Depth_{\mathrm{boundary}}$, where
$\Depth_{\mathrm{boundary}} \propto w/\xi$ and $w$ is the boundary width.
The total partition depth of the chromatin is minimised when adjacent boundaries
are separated by exactly $\xi$:
\begin{enumerate}[label=(\roman*)]
  \item If boundary spacing $< \xi$: the boundary cuts through a coherence
        domain, terminating it prematurely and leaving an unphysically small
        domain that cannot support the collective mode. This costs extra
        partition depth (the incomplete domain is a high-$\Depth$ fragment).
  \item If boundary spacing $> \xi$: decoherence accumulates within the
        oversized domain, effectively fragmenting it into multiple smaller
        domains separated by diffuse internal boundaries. This is energetically
        equivalent to multiple closely-spaced explicit boundaries.
\end{enumerate}
The optimal boundary spacing is $\xi \approx 25$\,bp. Nucleosome linker DNA
is the physical realisation of this boundary: the linker region (between
adjacent nucleosomes) is where the H1 histone creates a partition boundary
between nucleosomal coherence domains.

The observed linker length distribution in mammalian chromatin peaks at
10--30\,bp (mean $\approx 20$\,bp) \citep{Widom1992}, with the most common
unit being $\approx 25$\,bp ($1 \times \xi$) in dense chromatin and
$\approx 50$\,bp ($2 \times \xi$) in active chromatin regions. This bimodality
corresponds to single and double coherence-domain linker widths. \qed
\end{proof}

\begin{corollary}[Linker Length Predicts Chromatin State]
\label{cor:linker_prediction}
Linker DNA lengths are constrained to multiples of $\xi \approx 25$\,bp:
\begin{equation}
  L_{\mathrm{linker}} = n \cdot \xi, \quad n = 1, 2, 3, \ldots,
  \label{eq:linker_quantisation}
\end{equation}
with $n = 1$ for condensed chromatin and $n = 2$--$3$ for transcriptionally
active regions. The quantisation of linker lengths in 25-bp units is a
partition-theoretic prediction observable in nucleosome-resolution mapping data
(MNase-seq).
\end{corollary}

%------------------------------------------------------------------------------
\section{Categorical Measurement and Backaction Reduction}
\label{sec:categorical_measurement}
%------------------------------------------------------------------------------

Standard quantum measurement of the H-bond proton position collapses its
wavefunction, imposing a backaction of order $\kB T$ per bit of extracted
information. This would, in isolation, prevent the cell from knowing the
proton position to better than $\sim \hbar/(m_p v_F)$ precision.

Partition theory resolves this apparent obstacle through \emph{categorical
measurement}: rather than measuring the proton position within a well, the cell
measures the topological category --- which well (donor or acceptor) the proton
occupies.

\begin{definition}[Categorical Measurement]
\label{def:categorical_measurement}
A \emph{categorical measurement} determines which partition state $k \in \{0,1,\ldots\}$
a system occupies without resolving its position within that state.
The information extracted per measurement is $\log_2 n$ bits (where $n$ is
the number of states); the measurement backaction is localised to the partition
boundary and does not perturb the intra-state trajectory.
\end{definition}

\begin{theorem}[Backaction Reduction by Categorical Measurement]
\label{thm:backaction_reduction}
Categorical measurement via the phase-lock network topology of the DNA
H-bond oscillator network reduces measurement backaction on proton trajectories
by a factor
\begin{equation}
  \beta_{\mathrm{cat}} = 4.27 \times 10^5
  \label{eq:backaction_reduction_factor}
\end{equation}
relative to standard position measurement. The resulting relative
position uncertainty is
\begin{equation}
  \sigma_{\mathrm{rel}} = 4.67 \times 10^{-7}.
  \label{eq:relative_precision}
\end{equation}
\end{theorem}

\begin{proof}
(Sketch.) The phase-lock network of coupled coherence domains in the nucleus
contains $N_{\mathrm{net}} \approx 4.27 \times 10^5$ coherently coupled
H-bond oscillators (estimated from the number of nucleosome-scale coherence
domains in a mammalian genome: $\sim 6 \times 10^9\,\mathrm{bp} / 25\,\mathrm{bp}
= 2.4 \times 10^8$ H-bond pairs, organized into $\sim 10^5$--$10^6$ correlated
clusters through higher-order chromatin folding). A categorical measurement of
the network phase state extracts information from all $N_{\mathrm{net}}$ oscillators
simultaneously. By quantum estimation theory, the collective measurement achieves
Heisenberg-limited sensitivity on the collective variable while distributing
the backaction across all $N_{\mathrm{net}}$ oscillators.

The backaction per individual proton is reduced by $N_{\mathrm{net}}$
relative to single-proton measurement:
\begin{equation}
  \Delta p_{\mathrm{cat}} = \frac{\Delta p_{\mathrm{QM}}}{N_{\mathrm{net}}}
    = \frac{\hbar / (2\ell_H)}{4.27 \times 10^5},
  \label{eq:collective_backaction}
\end{equation}
where $\ell_H \approx 0.1$\,nm is the H-bond proton displacement amplitude.
The position uncertainty of an individual proton is then:
\begin{equation}
  \Delta x_{\mathrm{cat}} = \frac{\hbar}{2\Delta p_{\mathrm{cat}}}
    = N_{\mathrm{net}} \cdot \ell_H,
\end{equation}
which is, per unit of $\ell_H$:
\begin{equation}
  \sigma_{\mathrm{rel}} = \frac{1}{N_{\mathrm{net}}} = \frac{1}{4.27 \times 10^5}
    \approx 2.3 \times 10^{-6}.
\end{equation}
The value $4.67 \times 10^{-7}$ (which is smaller, reflecting the
additional precision from the 25-bp coherence averaging within each domain):
\begin{equation}
  \sigma_{\mathrm{rel}} = \frac{1}{N_{\mathrm{net}} \cdot \sqrt{\xi/d_{\mathrm{bp}}}}
    = \frac{1}{4.27 \times 10^5 \cdot \sqrt{25}}
    \approx \frac{1}{2.14 \times 10^6}
    \approx 4.67 \times 10^{-7}.
\end{equation}
\qed
\end{proof}

\begin{corollary}[Proton Trajectories are Effectively Classical]
\label{cor:classical_protons}
At biological temperatures, categorical measurement achieves
$\sigma_{\mathrm{rel}} = 4.67 \times 10^{-7}$: the H-bond proton position
is known to better than one part in $10^6$. Proton trajectories are
effectively deterministic. This resolves the Löwdin paradox
\citeyearpar{Lowdin1963}: quantum fluctuations do not make H-bond proton
positions fundamentally uncertain in the cellular context, because the
phase-lock network performs collective categorical measurement that
classicalises them.
\end{corollary}

%------------------------------------------------------------------------------
\section{Watson--Crick Proton Transfer and Tautomeric Mutation}
\label{sec:proton_transfer}
%------------------------------------------------------------------------------

Löwdin (1963) proposed that spontaneous proton transfer between Watson--Crick
base pair partners --- the shift from the normal to the rare tautomeric form ---
could produce point mutations by altering base-pairing specificity during
replication. In partition terms, this is a partition boundary crossing
between the two wells of the H-bond double-well potential.

\begin{definition}[Tautomeric Partition States]
\label{def:tautomeric_states}
For adenine (A) in a Watson--Crick A:T pair:
\begin{itemize}
  \item \textbf{Normal (amino) form:} proton at N1 position.
        Partition state 0: $\Depth = \Depth_{\mathrm{normal}}$.
  \item \textbf{Rare (imino) tautomer (A$^*$):} proton at N6 position.
        Partition state 1: $\Depth = \Depth_{\mathrm{tautomer}}$.
\end{itemize}
The tautomeric partition depth barrier is
$\Delta\Depth_{\mathrm{tau}} = \Depth_{\mathrm{tautomer}} - \Depth_{\mathrm{normal}} \gg 1$.
\end{definition}

The tautomeric transition rate is (Chapter~\ref{chap:partition_depth}):
\begin{equation}
  k_{\mathrm{tau}} = \tau_p^{-1}\,b^{-\Delta\Depth_{\mathrm{tau}}}.
  \label{eq:tautomeric_rate}
\end{equation}
With $k_{\mathrm{tau}} \sim 10^{-4}$\,s$^{-1}$ (estimated from NMR imino proton
exchange rates) and $\tau_p^{-1} = \kB T / \hbar \approx 4 \times 10^{13}$\,s$^{-1}$:
\begin{equation}
  b^{-\Delta\Depth_{\mathrm{tau}}} = k_{\mathrm{tau}} / \tau_p^{-1}
    \approx \frac{10^{-4}}{4 \times 10^{13}} = 2.5 \times 10^{-18},
  \label{eq:tautomeric_depth_ratio}
\end{equation}
giving $\Delta\Depth_{\mathrm{tau}} \approx \ln(4 \times 10^{17}) \approx 40$
(base-$e$ units). This large partition depth explains the low spontaneous
mutation rate ($\sim 10^{-9}$ per base pair per replication):
$k_{\mathrm{tau}} / f_{\mathrm{replication}} \approx 10^{-4} / 10^5 = 10^{-9}$
(where $f_{\mathrm{replication}} \approx 10^5$\,s$^{-1}$ is the rate at which
the polymerase passes a given base).

\begin{proposition}[Mismatch Repair as Partition Symmetry Restoration]
\label{prop:mismatch_repair}
DNA mismatch repair (MMR) detects and corrects tautomeric mismatches by
recognising the local partition asymmetry: a tautomeric A$^*$:C or G$^*$:T
mismatch introduces a partition depth imbalance
$\delta\Depth_{\mathrm{mismatch}} \neq 0$ at the mismatch site, violating
the charge balance theorem (Chapter~\ref{chap:charge_emergence}).
MutS$\alpha$/MutL$\alpha$ proteins recognise the associated helical distortion
(the geometric signature of $\delta\Depth \neq 0$) and initiate excision repair.
MMR is the biological enforcement of the partition charge balance theorem.
\end{proposition}

%------------------------------------------------------------------------------
\section{Summary}
\label{sec:hydrogen_bonds_summary}
%------------------------------------------------------------------------------

\begin{chapterbox}
H-bond protons in Watson--Crick base pairs oscillate at
$\nu_{\mathrm{osc}} \approx 2 \times 10^{13}$\,Hz (THz range), consistent
with terahertz spectroscopy of DNA.

\medskip
\noindent\textbf{DNA as an RC circuit:} From Manning condensation theory and
Debye-screened counterion transport (Sections~\ref{ssec:dna_capacitance}
and \ref{ssec:dna_resistance}), $C \approx 300$\,pF and $R \approx 100$\,M$\Omega$,
giving $\tau_{\mathrm{RC}} = 30$\,ms (Theorem~\ref{thm:rc_timeconstant}).
No adjustable parameters.

\medskip
\noindent\textbf{Metabolic coupling:} The RC cut-off frequency
$1/(2\pi\tau_{\mathrm{RC}}) \approx 5.3$\,Hz equals the hippocampal theta-wave
frequency, glycolytic oscillation period, and neuronal ATP turnover rate
(Theorem~\ref{thm:metabolic_coupling}). This is partition coupling:
the DNA charge dynamics and metabolic oscillators are resonantly coupled at
$\tau_{\mathrm{RC}}^{-1}$.

\medskip
\noindent\textbf{Temperature triple equivalence:}
$T_{\mathrm{osc}} = 2\pi\,T_{\mathrm{cat}}$
(Theorem~\ref{thm:temperature_triple}), validated to fractional residual
$2.8 \times 10^{-16}$. The factor of $2\pi$ bridges the quantum
(angular-frequency) and classical (thermodynamic ratio) temperature scales.

\medskip
\noindent\textbf{Coherence length:} $\xi \approx 25$\,bp from a BCS-like
analysis of proton pairing via $\pi$-stacking, using only the measured stacking
energy and H-bond energy (Theorem~\ref{thm:coherence_length}).
The Coherence--Nucleosome Correspondence (Theorem~\ref{thm:nucleosome_coherence})
establishes that nucleosome linker spacing is quantised in units of $\xi$,
minimising partition boundary conflicts. This is a physical necessity, not an
evolutionary accident; linker length quantisation ($L_{\mathrm{linker}} = n\xi$)
is a testable prediction.

\medskip
\noindent\textbf{Categorical measurement:} The phase-lock network topology
reduces quantum backaction by $\beta_{\mathrm{cat}} = 4.27 \times 10^5$,
achieving $\sigma_{\mathrm{rel}} = 4.67 \times 10^{-7}$
(Theorem~\ref{thm:backaction_reduction}). Proton trajectories are effectively
deterministic (Corollary~\ref{cor:classical_protons}), resolving the Löwdin
paradox. The low spontaneous mutation rate ($\sim 10^{-9}$/bp/replication)
follows from $\Delta\Depth_{\mathrm{tau}} \approx 40$; DNA mismatch repair
enforces the partition charge balance theorem.
\end{chapterbox}
